\documentclass[10pt,a4paper]{amsart}
\usepackage{amsmath,amssymb,amsthm,graphicx}
% Retained minimal preamble for the Stacks Project chapter more-algebra.tex (More on Algebra).
% The chapter file pinned-original-source.tex carries only \input{preamble} on its first line, so no
% preamble text lives inside the chapter itself.  The Stacks Project uses one single shared preamble for
% every chapter; the snapshot stages it identically as sources/stacks-chow/preamble.tex and
% sources/stacks-etale-cohomology/preamble.tex.  The block below is copied mechanically, line for line,
% from sources/stacks-chow/preamble.tex lines 14-18 (the Xy-pic setup) and 154-217 (the native theorem
% declarations together with the mathematical macros).  Nothing was added, removed or redefined; the
% wrapper class and the amsmath/amssymb/amsthm/graphicx base packages are supplied by the assembler.
% Chapter: more-algebra.tex -- More on Algebra.
% For commutative diagrams we use Xy-pic
\usepackage[all]{xy}

% We use 2cell for 2-commutative diagrams.
\xyoption{2cell}

% Theorem environments.
%
\theoremstyle{plain}
\newtheorem{theorem}[subsection]{Theorem}
\newtheorem{proposition}[subsection]{Proposition}
\newtheorem{lemma}[subsection]{Lemma}

\theoremstyle{definition}
\newtheorem{definition}[subsection]{Definition}
\newtheorem{example}[subsection]{Example}
\newtheorem{exercise}[subsection]{Exercise}
\newtheorem{situation}[subsection]{Situation}

\theoremstyle{remark}
\newtheorem{remark}[subsection]{Remark}
\newtheorem{remarks}[subsection]{Remarks}

\numberwithin{equation}{subsection}

% Macros
%
\def\lim{\mathop{\mathrm{lim}}\nolimits}
\def\colim{\mathop{\mathrm{colim}}\nolimits}
\def\Spec{\mathop{\mathrm{Spec}}}
\def\Hom{\mathop{\mathrm{Hom}}\nolimits}
\def\Ext{\mathop{\mathrm{Ext}}\nolimits}
\def\SheafHom{\mathop{\mathcal{H}\!\mathit{om}}\nolimits}
\def\SheafExt{\mathop{\mathcal{E}\!\mathit{xt}}\nolimits}
\def\Sch{\mathit{Sch}}
\def\Mor{\mathop{\mathrm{Mor}}\nolimits}
\def\Ob{\mathop{\mathrm{Ob}}\nolimits}
\def\Sh{\mathop{\mathit{Sh}}\nolimits}
\def\NL{\mathop{N\!L}\nolimits}
\def\CH{\mathop{\mathrm{CH}}\nolimits}
\def\proetale{{pro\text{-}\acute{e}tale}}
\def\etale{{\acute{e}tale}}
\def\QCoh{\mathit{QCoh}}
\def\Ker{\mathop{\mathrm{Ker}}}
\def\Im{\mathop{\mathrm{Im}}}
\def\Coker{\mathop{\mathrm{Coker}}}
\def\Coim{\mathop{\mathrm{Coim}}}

% Boxtimes
%
\DeclareMathSymbol{\boxtimes}{\mathbin}{AMSa}{"02}

%
% Macros for moduli stacks/spaces
%
\def\QCohstack{\mathcal{QC}\!\mathit{oh}}
\def\Cohstack{\mathcal{C}\!\mathit{oh}}
\def\Spacesstack{\mathcal{S}\!\mathit{paces}}
\def\Quotfunctor{\mathrm{Quot}}
\def\Hilbfunctor{\mathrm{Hilb}}
\def\Curvesstack{\mathcal{C}\!\mathit{urves}}
\def\Polarizedstack{\mathcal{P}\!\mathit{olarized}}
\def\Complexesstack{\mathcal{C}\!\mathit{omplexes}}
% \Pic is the operator that assigns to X its picard group, usage \Pic(X)
% \Picardstack_{X/B} denotes the Picard stack of X over B
% \Picardfunctor_{X/B} denotes the Picard functor of X over B
\def\Pic{\mathop{\mathrm{Pic}}\nolimits}
\def\Picardstack{\mathcal{P}\!\mathit{ic}}
\def\Picardfunctor{\mathrm{Pic}}
\def\Deformationcategory{\mathcal{D}\!\mathit{ef}}

% --- appended 2026-10-12 (Chief): the project own definitions of the reference/slogan/history
% display environments, copied mechanically from the shared Stacks preamble lines 8-12
% (\usepackage{verbatim} plus the three \newenvironment lines, which the chapter preamble keeps
% verbatim from the project preamble). Nothing added, removed or redefined: these environments
% carry no mathematical content and the project itself typesets them as comments. Without these
% three lines a counted statement whose own text uses them fails with Environment slogan undefined.
\usepackage{verbatim}
\newenvironment{reference}{\comment}{\endcomment}
\newenvironment{slogan}{\comment}{\endcomment}
\newenvironment{history}{\comment}{\endcomment}
\usepackage{cleveref}
\crefname{originallemma}{original lemma}{original lemmas}
\Crefname{originallemma}{Original lemma}{Original lemmas}
\makeatletter
\newlabel{more-algebra-lemma-henselization-P-ring}{{0C36}{source}{}{sourcepointer-1}{}}
\newlabel{more-algebra-lemma-henselization-P-ring@cref}{{[originallemma][0][]0C36}{[1][0][]source}{}{}{}}
\makeatother
\crefname{originallemma}{original lemma}{original lemmas}
\Crefname{originallemma}{Original lemma}{Original lemmas}
\makeatletter
\newlabel{lemma-henselization-pair-P-ring}{{0BKA}{source}{}{sourcepointer-2}{}}
\newlabel{lemma-henselization-pair-P-ring@cref}{{[originallemma][0][]0BKA}{[1][0][]source}{}{}{}}
\makeatother
\crefname{originallemma}{original lemma}{original lemmas}
\Crefname{originallemma}{Original lemma}{Original lemmas}
\makeatletter
\newlabel{lemma-check-P-ring-maximal-ideals}{{0BIU}{source}{}{sourcepointer-3}{}}
\newlabel{lemma-check-P-ring-maximal-ideals@cref}{{[originallemma][0][]0BIU}{[1][0][]source}{}{}{}}
\makeatother
\crefname{originallemma}{original lemma}{original lemmas}
\Crefname{originallemma}{Original lemma}{Original lemmas}
\makeatletter
\newlabel{lemma-fibres-henselization}{{07QQ}{source}{}{sourcepointer-4}{}}
\newlabel{lemma-fibres-henselization@cref}{{[originallemma][0][]07QQ}{[1][0][]source}{}{}{}}
\makeatother
\crefname{originallemma}{original lemma}{original lemmas}
\Crefname{originallemma}{Original lemma}{Original lemmas}
\makeatletter
\newlabel{lemma-henselization-noetherian}{{06LJ}{source}{}{sourcepointer-5}{}}
\newlabel{lemma-henselization-noetherian@cref}{{[originallemma][0][]06LJ}{[1][0][]source}{}{}{}}
\makeatother
\crefname{originallemma}{original lemma}{original lemmas}
\Crefname{originallemma}{Original lemma}{Original lemmas}
\makeatletter
\newlabel{proposition-fs-regular}{{07PM}{source}{}{sourcepointer-6}{}}
\newlabel{proposition-fs-regular@cref}{{[originallemma][0][]07PM}{[1][0][]source}{}{}{}}
\makeatother
\title{The henselization and strict henselization of a P-ring are P-rings}
\author{The Stacks Project Authors}
\date{stacks-project:more-algebra.tex:local-snapshot-20261012; GNU FDL 1.2 or later}
\begin{document}
\maketitle
\setcounter{subsection}{1}
\noindent\textit{Editorial note.} All delivered mathematical text is the author's own, copied from the pinned Stacks Project chapter file more-algebra.tex; no statement or proof step was rewritten, reordered or removed.\par\medskip
\noindent\textit{Editorial note.} The chapter file pinned-original-source.tex carries only \textbackslash{}input\{preamble\} on its first line, so no preamble text lives inside the chapter. The retained file sources/stacks-more-algebra/retained-preamble.tex is a mechanical, line-for-line copy of the project's single shared preamble as staged in the snapshot (Xy-pic setup plus theorem declarations and mathematical macros). Nothing was added, removed or redefined, and the excerpt serves only as preamble, never as a span coordinate.\par\medskip
\noindent\textit{Editorial note.} Both counted spans start at the \textbackslash{}begin line of their environment, so the statement and the author's proof are delivered whole; surrounding prose of the section is deliberately outside the counted spans and is carried as standard prerequisites instead.\par\medskip
\noindent\textit{Editorial note.} The citation fields (source\_id, source\_version, authors, source\_url, license.name, source\_title) are identical to those of every other card this lane has produced from this chapter, so the batch selector treats them as one source.\par\medskip
\noindent\textit{Allowed standard prerequisites (editorial).}\begin{itemize}
\item The henselian pair of a P-ring is a P-ring, and it suffices to check property P on formal fibres. Lemma \textbackslash{}ref\{lemma-henselization-pair-P-ring\} and Lemma \textbackslash{}ref\{lemma-check-P-ring-maximal-ideals\}, cited at proof lines 13162 and 13165-13166. Give the henselization case and reduce the strict henselization case to formal fibres.
\item The fibre of the strict henselization at a prime is a product of residue fields of the primes over the prime below. Lemma \textbackslash{}ref\{lemma-fibres-henselization\}, cited at proof lines 13171-13173. Turns the product of formal fibres into the single tensor product where the completion of the strict henselization appears.
\item The completion of the henselization is formally smooth over the completion of the ring, hence regular. Lemma \textbackslash{}ref\{lemma-henselization-noetherian\} and Proposition \textbackslash{}ref\{proposition-fs-regular\}, cited at proof lines 13180-13183. Makes the base change of property P along the completion available through condition (C).
\item Separability of the residue field extension of the strict henselization. condition (E) and Lemma \textbackslash{}ref\{lemma-fibres-henselization\}, used at proof lines 13193-13195. Transfers property P from the residue field of the prime below to the residue field of the prime.
\end{itemize}\medskip
\ifdefined\hypertarget\hypertarget{sourcepointer-1}{}\fi
\noindent\textit{Source reference.} Original lemma 0C36: focused-seeds/AGP/sources/stacks-snapshot-20261011/tags.txt, source lines 15181--15181. The lemma's own canonical tag, quoted so that readers can resolve the label inside the delivered text; it points at the counted statement itself and carries no proof content.\par\medskip
\ifdefined\hypertarget\hypertarget{sourcepointer-2}{}\fi
\noindent\textit{Source reference.} Original lemma 0BKA: focused-seeds/AGP/sources/stacks-snapshot-20261011/tags.txt, source lines 14553--14553. Navigation pointer to a statement the counted proof refers to but does not contain; nothing load-bearing is delegated to it.\par\medskip
\ifdefined\hypertarget\hypertarget{sourcepointer-3}{}\fi
\noindent\textit{Source reference.} Original lemma 0BIU: focused-seeds/AGP/sources/stacks-snapshot-20261011/tags.txt, source lines 14502--14502. Navigation pointer to a statement the counted proof refers to but does not contain; nothing load-bearing is delegated to it.\par\medskip
\ifdefined\hypertarget\hypertarget{sourcepointer-4}{}\fi
\noindent\textit{Source reference.} Original lemma 07QQ: focused-seeds/AGP/sources/stacks-snapshot-20261011/tags.txt, source lines 9840--9840. Navigation pointer to a statement the counted proof refers to but does not contain; nothing load-bearing is delegated to it.\par\medskip
\ifdefined\hypertarget\hypertarget{sourcepointer-5}{}\fi
\noindent\textit{Source reference.} Original lemma 06LJ: focused-seeds/AGP/sources/stacks-snapshot-20261011/tags.txt, source lines 8468--8468. Navigation pointer to a statement the counted proof refers to but does not contain; nothing load-bearing is delegated to it.\par\medskip
\ifdefined\hypertarget\hypertarget{sourcepointer-6}{}\fi
\noindent\textit{Source reference.} Original lemma 07PM: focused-seeds/AGP/sources/stacks-snapshot-20261011/tags.txt, source lines 9802--9802. Navigation pointer to a statement the counted proof refers to but does not contain; nothing load-bearing is delegated to it.\par\medskip
\begin{lemma}
\label{lemma-henselization-P-ring}
Let $R$ be a Noetherian local ring which is a $P$-ring where $P$
satisfies (B), (C), (D), and (E). Then the henselization $R^h$
and the strict henselization $R^{sh}$ are $P$-rings.
\end{lemma}

\begin{proof}
We have seen this for the henselization in
Lemma \ref{lemma-henselization-pair-P-ring}.
To prove it for the strict henselization, it suffices to
show that the formal fibres of $R^{sh}$ have $P$, see
Lemma \ref{lemma-check-P-ring-maximal-ideals}.
Let $\mathfrak r \subset R^{sh}$ be a prime and set
$\mathfrak p = R \cap \mathfrak r$. Set $\mathfrak r_1 = \mathfrak r$
and let $\mathfrak r_2, \ldots, \mathfrak r_s$
be the other primes of $R^{sh}$ lying over $\mathfrak p$, so that
$R^{sh} \otimes_R \kappa(\mathfrak p) =
\prod\nolimits_{i = 1, \ldots, s} \kappa(\mathfrak r_i)$, see
Lemma \ref{lemma-fibres-henselization}.
Then we see that
$$
\prod\nolimits_{i = 1, \ldots, t}
(R^{sh})^\wedge \otimes_{R^{sh}} \kappa(\mathfrak r_i) =
(R^{sh})^\wedge \otimes_{R^{sh}} (R^{sh} \otimes_R \kappa(\mathfrak p)) =
(R^{sh})^\wedge \otimes_R \kappa(\mathfrak p)
$$
Note that $R^\wedge \to (R^{sh})^\wedge$ is formally smooth
in the $\mathfrak m_{(R^{sh})^\wedge}$-adic topology, see
Lemma \ref{lemma-henselization-noetherian}.
Hence $R^\wedge \to (R^{sh})^\wedge$ is regular by
Proposition \ref{proposition-fs-regular}.
We conclude that property $P$ holds for
$\kappa(\mathfrak p) \to (R^{sh})^\wedge \otimes_R \kappa(\mathfrak p)$
by (C) and our assumption on $R$. Using property (B), using the
decomposition above, and looking
at local rings we conclude that property $P$ holds for
$\kappa(\mathfrak p) \to (R^{sh})^\wedge \otimes_{R^{sh}} \kappa(\mathfrak r)$.
Since $\kappa(\mathfrak r)/\kappa(\mathfrak p)$ is separable algebraic,
it follows from (E) that $P$ holds for
$\kappa(\mathfrak r) \to (R^{sh})^\wedge \otimes_{R^{sh}} \kappa(\mathfrak r)$.
\end{proof}

\end{document}
